\documentclass[11pt]{article}
\usepackage[margin=1in]{geometry}
\usepackage[T1]{fontenc}
\usepackage[utf8]{inputenc}
\usepackage{hyperref}
\usepackage{xcolor}
\usepackage{listings}
\usepackage{amsmath}
\usepackage{booktabs}
\usepackage{longtable}
\usepackage{enumitem}

\lstdefinestyle{code}{
  basicstyle=\ttfamily\small,
  breaklines=true,
  frame=single,
  columns=fullflexible,
  keepspaces=true,
  showstringspaces=false
}

\title{Llama 7B CUDA Kernel Handbook\\\large Modular, Kernel-wise Forward Pass Explanation}
\author{llama\_inference\_cpp}
\date{\today}

\begin{document}
\maketitle
\tableofcontents
\newpage

\section{Goal}
This handbook explains the modular CUDA kernel set written for Llama-7B-style inference in this repository:
\begin{itemize}[leftmargin=*]
  \item kernel by kernel,
  \item calculation by calculation,
  \item from the beginning to the end of the forward pass,
  \item with code split into readable files.
\end{itemize}

\section{Model context}
The 7B model used for this setup:
\begin{itemize}[leftmargin=*]
  \item checkpoint: \texttt{models/open\_llama\_7b\_v2.bin}
  \item tokenizer: \texttt{models/open\_llama\_7b\_v2\_tokenizer.bin}
\end{itemize}

Typical dimensions in this 7B setup:
\begin{itemize}[leftmargin=*]
  \item \(\texttt{dim}=4096\)
  \item \(\texttt{n\_layers}=32\)
  \item \(\texttt{n\_heads}=32\)
  \item \(\texttt{head\_size}=128\)
  \item \(\texttt{hidden\_dim}=11008\)
\end{itemize}

\section{Modular kernel file layout}
\begin{longtable}{p{0.43\linewidth} p{0.5\linewidth}}
\toprule
\textbf{File} & \textbf{Purpose} \\
\midrule
\endhead
\texttt{include/cuda\_kernels/llama\_cuda\_kernels.cuh} & Kernel launch API (one function per op). \\
\texttt{src/cuda\_kernels/rmsnorm\_kernel.cu} & RMSNorm computation kernel. \\
\texttt{src/cuda\_kernels/matvec\_kernel.cu} & Matrix-vector projection kernel. \\
\texttt{src/cuda\_kernels/rope\_kernel.cu} & RoPE position rotation for q and k. \\
\texttt{src/cuda\_kernels/attention\_scores\_kernel.cu} & Dot-product attention scores per timestep. \\
\texttt{src/cuda\_kernels/softmax\_kernel.cu} & In-place stable softmax for score vectors. \\
\texttt{src/cuda\_kernels/attention\_weighted\_sum\_kernel.cu} & Value weighted-sum per head. \\
\texttt{src/cuda\_kernels/swiglu\_kernel.cu} & SwiGLU elementwise nonlinearity. \\
\bottomrule
\end{longtable}

\section{Forward pass from start to end (kernel order)}
For one token at decoding position \texttt{pos}:

\begin{enumerate}[leftmargin=*]
  \item Embedding lookup: \(x \leftarrow E[token]\).
  \item For each layer \(l\):
  \begin{enumerate}
    \item \texttt{launch\_rmsnorm} on \(x\) for attention branch.
    \item \texttt{launch\_matvec} for \(W_qx, W_kx, W_vx\).
    \item \texttt{launch\_rope\_qk} on \(q, k\).
    \item Write \(k,v\) into KV cache at position \texttt{pos}.
    \item For each head:
      \begin{itemize}
        \item \texttt{launch\_attention\_scores}
        \item \texttt{launch\_softmax\_inplace}
        \item \texttt{launch\_attention\_weighted\_sum}
      \end{itemize}
    \item \texttt{launch\_matvec} for \(W_o\) and attention residual add.
    \item \texttt{launch\_rmsnorm} for FFN branch.
    \item \texttt{launch\_matvec} for \(W_1x\) and \(W_3x\).
    \item \texttt{launch\_swiglu}.
    \item \texttt{launch\_matvec} for \(W_2\) and FFN residual add.
  \end{enumerate}
  \item Final \texttt{launch\_rmsnorm}.
  \item \texttt{launch\_matvec} with \(W_{cls}\) to produce logits.
\end{enumerate}

\section{Kernel-by-kernel calculations}
\subsection{RMSNorm kernel}
\[
\text{ss} = \frac{1}{\sqrt{\frac{1}{d}\sum_{j=0}^{d-1}x_j^2 + \epsilon}},
\quad
out_j = weight_j \cdot (x_j \cdot \text{ss})
\]

Code:
\lstinputlisting[style=code]{../src/cuda_kernels/rmsnorm_kernel.cu}

\subsection{MatVec kernel}
\[
out_i = \sum_{j=0}^{n-1} W_{i,j} x_j
\]

Code:
\lstinputlisting[style=code]{../src/cuda_kernels/matvec_kernel.cu}

\subsection{RoPE kernel}
For channel pair \((x_0,x_1)\):
\[
\begin{bmatrix}x_0'\\x_1'\end{bmatrix}
=
\begin{bmatrix}\cos\theta & -\sin\theta\\ \sin\theta & \cos\theta\end{bmatrix}
\begin{bmatrix}x_0\\x_1\end{bmatrix}
\]

Code:
\lstinputlisting[style=code]{../src/cuda_kernels/rope_kernel.cu}

\subsection{Attention score kernel}
\[
score_t = \frac{q\_{head}\cdot k\_{head,t}}{\sqrt{head\_size}}
\]

Code:
\lstinputlisting[style=code]{../src/cuda_kernels/attention_scores_kernel.cu}

\subsection{Softmax kernel}
\[
p_t = \frac{\exp(score_t - \max(score))}{\sum_u \exp(score_u - \max(score))}
\]

Code:
\lstinputlisting[style=code]{../src/cuda_kernels/softmax_kernel.cu}

\subsection{Attention weighted-sum kernel}
\[
out_i = \sum_{t=0}^{T-1} p_t \cdot v_{t,i}
\]

Code:
\lstinputlisting[style=code]{../src/cuda_kernels/attention_weighted_sum_kernel.cu}

\subsection{SwiGLU kernel}
\[
\mathrm{SiLU}(x)=x\cdot\sigma(x),\quad hb_i = \mathrm{SiLU}(hb_i)\cdot hb2_i
\]

Code:
\lstinputlisting[style=code]{../src/cuda_kernels/swiglu_kernel.cu}

\section{Single API header for readability}
Public launch interfaces (one function per kernel):
\lstinputlisting[style=code]{../include/cuda_kernels/llama_cuda_kernels.cuh}

\section{Build}
\begin{lstlisting}[style=code]
cmake -S . -B build -DCMAKE_BUILD_TYPE=Release -DLLAMA_INFERENCE_BUILD_CUDA_KERNELS=ON
cmake --build build -j
\end{lstlisting}

\section{Why this is modular and readable}
\begin{itemize}[leftmargin=*]
  \item Each CUDA operation is isolated in its own file.
  \item One shared API header defines the launch surface.
  \item Forward-pass order maps directly to the kernel sequence.
  \item You can profile and optimize each kernel independently.
\end{itemize}

\end{document}
